\documentclass[12pt]{article}

\usepackage[utf8]{inputenc}
\usepackage[french]{babel}
\usepackage[T1]{fontenc}
\usepackage{geometry}
\geometry{margin=2cm}
\usepackage{listings}
\usepackage{xcolor}
\usepackage{csquotes}

% --- Packages ---
\usepackage[utf8]{inputenc}
\usepackage[T1]{fontenc}
\usepackage[french]{babel}
\usepackage{amsmath, amssymb}
\usepackage{geometry}
% \usepackage[draft]{graphicx}
\usepackage{tikz}
\usepackage[most]{tcolorbox}
\usepackage{multicol}
\usepackage{lmodern}
\usepackage{pgfplots}
\pgfplotsset{compat=1.18}
\usepackage{mdframed}
\usepackage{array}
\usepackage[thinlines]{easytable}
\usepackage{tkz-tab}
\usepackage{bm}
\usepackage{xcolor}
\usepackage{sectsty}
\usepackage[utf8]{inputenc}
\usepackage[T1]{fontenc}
\usepackage[inline]{enumitem}
\usepackage[dvipsnames]{xcolor}
\usepackage{multido}
% --- Page ---
\geometry{margin=2cm}

% Footer
\usepackage{fancyhdr}
\pagestyle{fancy}
\cfoot{\thepage}
\rfoot{\textit{G.M.}}
\renewcommand{\headrulewidth}{0pt}
\renewcommand{\footrulewidth}{0pt}

% Pour les normes : 
\newcommand*{\norme}[1]{\left\lVert\overrightarrow{#1}\right\rVert}

\usetikzlibrary{shapes.misc}

\tikzset{cross/.style={cross out, draw=black, minimum size=2*(#1-\pgflinewidth), inner sep=0pt, outer sep=0pt},
%default radius will be 1pt. 
cross/.default={4pt}}

% mettre dans le préambule, après \usepackage{xcolor}
\usepackage{sectsty} % gère facilement la couleur des titres
\setlength{\parindent}{0pt}

\sectionfont{\color{red}\Large\bfseries}     % pour \section*
\subsectionfont{\color{red}\large\bfseries\hspace{1cm}}

% Configuration pour le code Python
\lstset{
    language=Python,
    basicstyle=\ttfamily\footnotesize,
    keywordstyle=\color{blue}\bfseries,
    commentstyle=\color{green!50!black},
    stringstyle=\color{red},
    showstringspaces=false,
    breaklines=true,
    frame=single,
    numbers=left,
    numberstyle=\tiny\color{gray},
    captionpos=b,
    xleftmargin=15pt,
    framesep=5pt
}

% \title{Fiche Recapitulative - Python (Seconde)}
% \author{}
% \date{}

\begin{document}

\begin{center}
    {\LARGE \textbf{Python -- Algorithmique}}\\[4mm]
    {\large - 2\up{nde} -}
\end{center}

\vspace{10pt}
\hrule
\vspace{0pt}

\vspace{10pt}
\section*{\textcolor{red}{I/ Les fonctions}}
Une \textbf{fonction} est un bloc de code reutilisable, defini avec le mot-cle \texttt{def}.

\subsection*{\hspace{1cm} 1) Définition et appel}
\vspace{5pt}

\begin{lstlisting}[caption={Definition et utilisation}]
# Definition
def dire_bonjour(nom):
    print(f"Bonjour, {nom}!")

# Appel
dire_bonjour("Alice")  # Affiche "Bonjour, Alice!"
\end{lstlisting}

\subsection*{\hspace{1cm} 2) Valeur de retour}
\vspace{5pt}
\begin{lstlisting}[caption={Fonction avec retour}]
def carre(x):
    return x * x

resultat = carre(4)  # resultat = 16
print(resultat)
\end{lstlisting}

\begin{mdframed}[
    linewidth=1pt,
    roundcorner=5pt,
    linecolor=RoyalBlue,
    innertopmargin=10pt,
    innerbottommargin=10pt
]
{\color{RoyalBlue}
\textbf{\underline{Exercice d'application n°1:} Utiliser des fonctions.}}

\vspace{5pt}
Pour définir avec Python la fonction $f$ telle que $f(x) = 2x+3$, il faut saisir :

\begin{minipage}[t]{.45\textwidth}%
    \vspace{0pt}
    \begin{lstlisting}
    def f(x):
        return 2*x+3
    \end{lstlisting}
\end{minipage}%
\hfill
\begin{minipage}[t]{.45\textwidth}%
    \vspace{0pt}
    
\end{minipage}%

\hspace{1cm}1) Modifier le programme pour que la fonction retourne le résultat de 

\hspace{1.4cm} $g(x) = -2x^2+3x-7$.

\vspace{0.5cm}
{\color{black!45}\dotfill}

\vspace{0.5cm}
{\color{black!45}\dotfill}

\vspace{0.5cm}
{\color{black!45}\dotfill}

\vspace{0.5cm}
{\color{black!45}\dotfill}


\newpage
{\color{RoyalBlue}
\textbf{\underline{Exercice d'application n°2:}}}

\vspace{5pt}
Écrire un programme affichant le prix réduit en fonction du
prix de départ et de la réduction accordée en \%.

\begin{lstlisting}
    .
    .
    .
    .
    .
    .
    .
    .
    .
    .

\end{lstlisting}

\vspace{20pt}

\end{mdframed}


\section*{\textcolor{red}{II/ L'affectation}}
L'\textbf{affectation} permet de stocker une valeur dans une variable a l'aide du symbole \texttt{=}.

\vspace{5pt}
\begin{lstlisting}[caption={Exemples d'affectation}]
# Affectation simple
x = 5
nom = "Alice"
est_majeure = False # ce type de variable est appelee "booleen"

# Affectation multiple
a, b = 10, 20

# Echange de valeurs
a, b = b, a
\end{lstlisting}

\vspace{10pt}
\section*{\textcolor{red}{III/ Conditions}}
Les conditions permettent d'executer du code en fonction de tests logiques.

\subsection*{1) Structure de base}
\vspace{5pt}

\begin{lstlisting}[caption={Structure \texttt{if-elif-else}}]
if condition1: # la condition est de type booleen
    # Code execute si condition1 est vraie
elif condition2:
    # Code execute si condition2 est vraie
else:
    # Code execute si aucune condition n'est vraie
\end{lstlisting}

\subsection*{2) Operateurs de comparaison}
\vspace{5pt}

\begin{center}
\begin{tabular}{|c|c|}
    \hline
    \textbf{Operateur} & \textbf{Signification} \\
    \hline
    \texttt{==} & egal a \\
    \hline
    \texttt{!=} & different de \\
    \hline
    \texttt{>} & superieur a \\
    \hline
    \texttt{>=} & superieur ou egal a \\
    \hline
    \texttt{<} & inferieur a \\
    \hline
    \texttt{<=} & inferieur ou egal a \\
    \hline
\end{tabular}
\end{center}

\subsection*{3) Operateurs logiques}
\begin{lstlisting}[caption={Combinaison de conditions}]
if (age >= 18) and (age <= 60):
    print("Adulte")
if (note < 10) or (absences > 5):
    print("Attention !")
\end{lstlisting}

\begin{mdframed}[
    linewidth=1pt,
    roundcorner=5pt,
    linecolor=RoyalBlue,
    innertopmargin=10pt,
    innerbottommargin=10pt
]
{\color{RoyalBlue}
\textbf{\underline{Exercice d'application n°3:} Tester des conditions en Python.}}

\vspace{5pt}
Dans le programme ci-dessous, les affichages en sortie de l'algorithme ont été supprimées.

\begin{lstlisting}
def signe(a,b):
    if a>0:
        if b>0:
            print("le produit axb est .....................")
        else:
            print("le produit axb est .....................")
    else:
        if b>0:
            print("le produit axb est .....................")
        else:
            print("le produit axb est .....................")
\end{lstlisting}

\vspace{5pt}
\hspace{1cm}1) Quel problème permet de résoudre cet
algorithme ?

\hspace{1cm}2) Compléter le programme par les
affichages en sortie manquants.

\vspace{0.5cm}
{\color{black!45}\dotfill}

\vspace{0.5cm}
{\color{black!45}\dotfill}

\vspace{0.5cm}
{\color{black!45}\dotfill}

\vspace{0.5cm}
{\color{black!45}\dotfill}


\vspace{20pt}

\newpage
{\color{RoyalBlue}
\textbf{\underline{Exercice d'application n°4:} Tester des conditions en Python.}}

\vspace{5pt}
Écrire et tester un programme qui demande en entrée à un client le montant total de
ses achats.

\vspace{5pt}
En fonction de la somme dépensée, le programme affiche en sortie le prix à payer :
\begin{itemize}
    \item Si la somme dépensée est strictement inférieure à 75 €, il obtient 5\% de remise.
    \item Si la somme dépensée est supérieure à 75 €, il obtient 8\% de remise.
\end{itemize}

\begin{lstlisting}
    .
    .
    .
    .
    .
    .
    .
    .
    .
    .
    .
    .
    .
    .
    .
    .
    .

\end{lstlisting}

\end{mdframed}

\vspace{10pt}
\section*{\textcolor{red}{IV/ Boucles}}
\subsection*{1) Boucle \texttt{for}}
La boucle \texttt{for} permet de repeter une action pour chaque element d'une sequence.
\begin{lstlisting}[caption={Exemples avec \texttt{for}}]
# Parcourir une liste
for element in [1, 2, 3]:
    print(element)

# Utilisation de range()
for i in range(5):  # 0 a 4
    print(i)

for i in range(2, 6):  # 2 a 5
    print(i)

for i in range(0, 10, 2):  # 0, 2, 4, 6, 8
    print(i)
\end{lstlisting}

\subsection*{2) Boucle \texttt{while}}
La boucle \texttt{while} repete une action \textbf{tant qu'}une condition est vraie.
\begin{lstlisting}[caption={Exemple avec \texttt{while}}]
compteur = 0
while compteur < 5:
    print(compteur)
    compteur += 1  # Incrementation
\end{lstlisting}

\begin{mdframed}[
    linewidth=1pt,
    roundcorner=5pt,
    linecolor=RoyalBlue,
    innertopmargin=10pt,
    innerbottommargin=10pt
]
{\color{RoyalBlue}
\textbf{\underline{Exercice d'application n°5:} Tester les boucles en Python.}}

\vspace{5pt}
\begin{minipage}[t]{.45\textwidth}%
    \vspace{0pt}
    1) Tester le programme ci-contre. Qu’affiche-t-il en sortie ?

    \vspace{0.5cm}
    {\color{black!45}\dotfill}
    
\end{minipage}%
\hfill
\begin{minipage}[t]{.5\textwidth}%
    \vspace{0pt}
\begin{lstlisting}
n=0
while n<10:
    print(n)
    n=n+2
\end{lstlisting}
\end{minipage}%

\vspace{0.5cm}
{\color{black!45}\dotfill}

\vspace{0.5cm}
{\color{black!45}\dotfill}

\vspace{20pt}

2) Écrire un programme qui affiche tous les entiers pairs
compris entre 18 et 45.

\begin{lstlisting}
.









.
\end{lstlisting}

\vspace{5pt}
3) Ecrire et tester un programme qui affiche tous les entiers
impairs
compris entre 50 et 150.

\begin{lstlisting}
.









.
\end{lstlisting}

\newpage
{\color{RoyalBlue}
\textbf{\underline{Exercice d'application n°5:} Tester les boucles en Python.}}

\vspace{5pt}
On place un capital de 500€ sur un compte rémunéré
à 3\% par an.

L’algorithme ci-contre, écrit en langage naturel,
permet de calculer le nombre d'années au bout
desquelles le capital sera doublé.

\begin{minipage}[t]{.45\textwidth}%
    \vspace{0pt}
    \hspace{1cm} 1) Le programme ci-contre traduisant l'algorithme précédent
    comprend une erreur. Corriger le programme et le tester.

    \vspace{5pt}
    \hspace{1cm} 2) Modifier le programme précédent de telle sorte que le capital et
    le taux de rémunération soient saisis en entrée.
\end{minipage}%
\hfill
\begin{minipage}[t]{.5\textwidth}%
    \vspace{0pt}
\begin{lstlisting}
def capital():
    s = 500
    a = 0
    while s<1000:
        s = 1.03*s
    a=a+1
    print(a)
\end{lstlisting}
\end{minipage}%
\vspace{5pt}


\begin{lstlisting}
.



















.
\end{lstlisting}

\end{mdframed}

\end{document}